\documentclass[10pt,a4paper]{amsart}
\usepackage[margin=19mm]{geometry}
\usepackage{amsmath,amssymb,mathtools}
\usepackage[scr=rsfs,cal=euler]{mathalfa}
\usepackage{xfrac}
\usepackage[unicode,hidelinks,bookmarks=false]{hyperref}
\newcommand{\ie}{i.e.\ }
\newcommand{\cf}{cf.\ }
\newcommand{\resp}{respectively\ }
\newcommand{\Subsection}{\subsection}
\providecommand{\nobreakdash}{\nobreak\hbox{-}}
\setlength{\emergencystretch}{2em}
\multlinegap0pt
\usepackage{amssymb}
\usepackage{xcolor}
\newtheorem{proposition}{Proposition}
\title{Minimal Two-Spheres and Manifolds with Positive Isotropic Curvature}
\author{Tsz-Kiu Aaron Chow, Yipeng Wang}
\date{Source: arXiv:2609.06910v1 (CC BY 4.0)}
\begin{document}
\maketitle

\noindent\emph{Source and licence.} Minimal Two-Spheres and Manifolds with Positive Isotropic Curvature. Version arXiv:2609.06910v1 (CC BY 4.0), distributed by arXiv under the Creative Commons Attribution 4.0 International licence, http://creativecommons.org/licenses/by/4.0/, as recorded in the arXiv licence metadata for this exact version and shown on the abstract page https://arxiv.org/abs/2609.06910v1. Full licence notice: \texttt{licenses/arxiv-2609.06910-CC-BY-4.0.txt}.\par\smallskip

\noindent\textit{Editorial scope.} Counted unit: the article's proposition prop:intermediate (source0.tex lines 115--117). The article's complete original proof of it is delivered with it (source0.tex lines 118--155, one \texttt{proof} environment of 38 lines). No mathematics is written, repaired or re-derived: the statement and the proof are the article's own lines, in the article's own notation, and no retained line is substituted or re-flowed.  No further source-local context is needed: every cross-reference of every retained line resolves inside the two delivered spans.  Nothing was omitted from any retained span, so the package carries no omission record.  No retained line contains a citation, so the wrapper carries no bibliography and no bibliography entry.  No retained line contains a figure environment, an image-inclusion command or drawing code, so no image is delivered and the package declares no asset dependency.  Editorial formatting only, in addition: an AMS \texttt{amsart} preamble that restates the article's own declarations and macros used by the retained lines, namely the environments \texttt{proposition} on separate counters, together with the standard packages those lines need.  The retained environments print in this document as Proposition 1 in order of appearance, whereas the article prints the same items with the numbering of its own full document.  The complete native source of the article as distributed in the arXiv e-print for this exact version is bundled unchanged as \texttt{sources/209589/source0.tex}.
\begin{proposition}\label{prop:intermediate}
Let $(M,g)$ be a Riemannian manifold of dimension $n\geq 4$ with positive isotropic curvature. Then $(M,g)$ has positive $m$-intermediate curvature for every $m\in \{3,\cdots, n-1\}$.     
\end{proposition}

\begin{proof}
    Let us fix some $p\in M$ and an orthonormal frame $\{e_1,\cdots,e_n\}$ at $p$. For every four distinct indices $i,j,k,\ell$, applying the positive isotropic curvature condition to the four-frames $\{e_i,e_j,e_k,e_{\ell}\}$ and $\{e_i,e_j,e_k,-e_{\ell}\}$ gives
    \begin{equation}\label{eqn:four-curvatures}
    R_{ikik}+R_{i\ell i\ell}+R_{jkjk}+R_{j\ell j\ell}>0.
    \end{equation}
    Clearly, $(M,g)$ has positive scalar curvature by summing over \eqref{eqn:four-curvatures} over every orthonormal four-frames. Hence, we fix some positive integer $m\in \{3,\cdots,n-2\}$. We now consider two cases:

    \textit{Case 1.} Suppose $4\leq m\leq n-2$. Then for every $1\leq i<j<k<\ell\leq m$, applying \eqref{eqn:four-curvatures} to the indices $(i,j,k,\ell)$, $(i,k,j,\ell)$ and $(i,\ell,j,k)$ gives
    \[
    2(R_{ijij}+R_{ikik}+R_{i\ell i\ell}+R_{jkjk}+R_{j\ell j\ell}+R_{k\ell k\ell})>0.
    \]
    Summing over every such $1\leq i<j<k<\ell\leq m$ gives
    \[
        (m-2)(m-3)\sum_{1\leq i<j\leq m}R_{ijij}>0.
    \]
    Moreover, for every $1\leq i<j\leq m$ and $m<k<\ell\leq n$, summing over \eqref{eqn:four-curvatures} for the indices $\{i,j,k,\ell\}$ gives
    \[
        (m-1)(n-m-1)\sum_{1\leq i\leq m,\,m<k\leq n}R_{ikik}>0.
    \]
    The claim then follows from the fact that  
    \[\sum_{i=1}^m\sum_{j=i+1}^nR_{ijij}=\sum_{1\leq i<j\leq m}R_{ijij}+\sum_{1\leq i\leq m,\,m<k\leq n}R_{ikik}.
    \]
    
    \textit{Case 2.} Suppose $m=3$. Then for every $k\in \{4,\cdots,n\}$, applying \eqref{eqn:four-curvatures} to the indices $(1,2,3,k)$, $(1,3,2,k)$ and $(1,k,2,3)$ gives
    \[
    2(R_{1212}+R_{1313}+R_{2323}+R_{1k1k}+R_{2k2k}+R_{3k3k})>0.
    \]
    Summing over $4\leq k\leq n$ gives $$2(n-3)(R_{1212}+R_{1313}+R_{2323})+2\sum_{k=4}^n(R_{1k1k}+R_{2k2k}+R_{3k3k})>0.$$ 
    Moreover, adding \eqref{eqn:four-curvatures} for all $1\leq i<j\leq 3$ and $4\leq k<\ell\leq n$ gives
    \[
    2(n-4)\sum_{k\geq 4}(R_{1k1k}+R_{2k2k}+R_{3k3k})>0.
    \]
    Adding these inequalities gives
    \[
    2(n-3)(R_{1212}+R_{2323}+R_{1313}+\sum_{i=1}^3\sum_{k=4}^nR_{ikik})>0.
    \]
    Finally, since $$\sum_{i=1}^3\sum_{j=i+1}^nR_{ijij}=R_{1212}+R_{1313}+R_{2323}+\sum_{i=1}^3\sum_{k=4}^nR_{ikik},$$ the claim follows directly.
\end{proof}

\end{document}
